\section{Theorem 1 in the life cycle: how far it goes}\label{sec:thm1}

The programme of Section~\ref{sec:prog} rests on Light's Theorem 1, that saving rises with the
interest rate at every state. In the infinite horizon this is open above logarithmic utility. The
life cycle was the reason for hope: a finite backward induction from an explicit terminal
condition, with the per-step loss compounding only $J$ times instead of to a fixed point. This
section reports what the numbers say, at $\beta=0.96$, a seven-state Tauchen chain with
$\rho=0.9$ and $\sigma=0.4$, a zero borrowing limit, $J=60$, and a flat earnings profile
(\texttt{numerics/lc\_check.m}, \texttt{lc\_robust.m}, \texttt{lc\_induct.m},
\texttt{lc\_verify.m}, \texttt{lc\_bound.m}, \texttt{lc\_crit.m}).

\subsection{The statement is its own terminal condition}

The budget constraint makes two statements identical. Since $a'=Ra+y(z)-c$ at either rate,
\[
  g_k(a,z;r_2)\ \ge\ g_k(a,z;r_1)
  \qquad\Longleftrightarrow\qquad
  c_k(a,z;r_2)-c_k(a,z;r_1)\ \le\ (r_2-r_1)\,a ,
\]
so Theorem 1 \emph{is} the affine bound on the rate response of consumption with slope one and
zero intercept. That bound holds with equality in the last period of life, where $c_0=y(z)+Ra$.
The question is therefore whether the terminal condition propagates backward, and it is a
question about one inequality rather than about a pair of constants to be chosen.

\subsection{It propagates, up to a critical risk aversion}

At $\mu=1$ and $\mu=3$ it propagates exactly. Over all sixty stages and the whole wealth region
reachable from zero wealth at birth, the largest violation of $g_k^{(1)}\le g_k^{(2)}$ is zero,
and the consumption bound holds to five parts in $10^{14}$; refining the asset grid from $2000$
to $16000$ points changes nothing. At $\mu=5$ it fails. The failure is a genuine feature of the
household's problem, not of the solution method: the violation is identical at every grid size,
and endogenous-grid and value-function iteration with a pure grid search agree on it to the
resolution of the grid. Saving at zero wealth in the highest income state, in the first period of
a sixty-period life, falls as the rate rises:

\begingroup\small
\begin{center}
\begin{tabular}{@{}lccccc@{}}
\toprule
$r$ & $2\%$ & $3\%$ & $4\%$ & $5\%$ & $6\%$ \\
\midrule
$g(0,z_{\max})$, endogenous grid & $1.7783$ & $1.7638$ & $1.7452$ & $1.7234$ & $1.6993$ \\
$g(0,z_{\max})$, value iteration & $1.7806$ & $1.7606$ & $1.7406$ & $1.7206$ & $1.7006$ \\
\bottomrule
\end{tabular}
\end{center}
\endgroup

Between the two lies a critical degree of risk aversion, and it is sharp. Bisecting on $\mu$ for
the largest violation over all stages gives:

\begingroup\small
\begin{center}
\begin{tabular}{@{}lccc@{}}
\toprule
Horizon $J$ & $r=2\%$ & $r=4\%$ & $r=\lambda=4.17\%$ \\
\midrule
$20$ & $3.79$ & $3.53$ & $3.51$ \\
$40$ & $3.77$ & $3.39$ & $3.36$ \\
$60$ & $3.77$ & $3.39$ & $3.36$ \\
\bottomrule
\end{tabular}
\end{center}
\endgroup

Theorem 1 holds throughout the life cycle below the critical value and fails above it. The value
is stable in the horizon beyond about forty periods, and falls as the rate rises towards the rate
of time preference. Aiyagari's $\mu=3$ is inside it with little room, at $\mu=3.39$ against
$3$; his $\mu=5$ is well outside. So the life cycle buys one of the two open cases and not the
other, and it buys $\mu=3$ only at rates below the rate of time preference, which is where the
equilibrium lies.

\subsection{Why the natural induction does not close}

The backward induction that the transfer lemma suggests compares the two Euler equations at a
state where the lower-rate household saves. Writing $b=g^{(1)}_{k+1}(a,z)$ for the pivot, the
induction hypothesis and monotonicity of $c_k$ give
$c^{(2)}_k(b')\le c^{(1)}_k(b)+\Delta r\,b$ at the relevant points, and the step closes provided
\[
  R_2\,\E\bigl[(c^{(1)}_k(b)+\Delta r\,b)^{-\mu}\bigr]\ \ge\
  R_1\,\E\bigl[c^{(1)}_k(b)^{-\mu}\bigr]\Bigl(1+\frac{\Delta r\,a}{c^{(1)}_{k+1}(a,z)}\Bigr)^{-\mu}.
\]
To first order this asks $\mu\bigl(b/c^{(1)}_k(b)-a/c^{(1)}_{k+1}(a,z)\bigr)\le1/R$: the ratio of
assets to consumption may rise from one age to the next by at most $1/(\mu R)$. It does not. At
zero wealth in the highest income state the left-hand side is large, because such a household
saves a great deal against a certain fall in income, while the right-hand side is about one. The
margin is negative from the third stage of life onward for every degree of risk aversion we tried,
logarithmic utility included, at $-0.12$ for $\mu=1$, $-2.86$ for $\mu=3$ and $-6.29$ for $\mu=5$
per unit of the rate gap. Keeping the term the step discards, the marginal propensity to consume
times the size of the hypothetical violation, bounds the violation by a multiple of $\Delta r$
rather than eliminating it, which is not enough: the aggregate difference the equilibrium argument
must sign is itself of order $\Delta r$.

\subsection{The reduction, formalised, and its verdict}

The backward induction above is proved in Lean, in the form the life cycle makes natural.
Because the budget constraint makes the two statements identical
(\lean{stagePolicy\_le\_iff\_consumption\_le}), the terminal stage needs no hypothesis at all: in
the last period both households consume their cash on hand and save nothing
(\lean{stagePolicy\_mono\_zero}). Corners need none either: if the lower-rate household saves
nothing there is nothing to violate. What the step needs is one inequality per age and state,
which we call the \emph{age-indexed slack condition}: writing $b$ for the lower-rate household's
saving at that state and $C'(z')$ for its consumption next period,
\begin{equation}\label{eq:stageslack}
  R_1\sum_{z'}\pi(z,z')C'(z')^{-\mu}\ \le\
  R_2\sum_{z'}\pi(z,z')\bigl(C'(z')+\Delta r\,b\bigr)^{-\mu}
  \Bigl(1+\frac{\Delta r\,a}{c}\Bigr)^{\mu}
\end{equation}
(\lean{StageSlack}).

\begin{theorem}[The reduction]\label{thm:reduction}
If \eqref{eq:stageslack} holds at every age and every state in the region, then saving rises with
the rate at every age, and aggregate capital rises with the rate.
\formal{IncomeFluctuation.stagePolicy\_mono\_step,
IncomeFluctuation.stagePolicy\_mono\_of\_stageSlack,
IncomeFluctuation.olgCapital\_le\_of\_stageSlack}
\end{theorem}

With Theorem~\ref{thm:thm3} this closes the chain: \eqref{eq:stageslack} at every age gives a
unique equilibrium rate. The question is therefore whether \eqref{eq:stageslack} is true, and the
answer is the one reported above, now stated sharply. Condition \eqref{eq:stageslack} is exactly
the margin tabulated in the previous subsection. At $\mu=3$ it holds for the last three or four
periods of life, with margins $+0.68$, $+0.42$ and $+0.18$ per unit of the rate gap, and fails
from about the fifth-to-last, at $-0.23$, deteriorating to $-2.86$ at birth. At $\mu=1$ it holds
for most of the life and still fails for the youngest ages, at $-0.12$ at birth.

That last observation is the important one. The condition fails even at logarithmic utility,
where Theorem 1 is certainly true. So the failure is not a symptom of high risk aversion: the
pointwise Euler comparison is simply the wrong instrument in a life cycle, and no sharpening of
it that we could find repairs the young ages. Keeping the term the step discards, the marginal
propensity to consume times the size of the hypothetical violation, only bounds the violation by
a multiple of $\Delta r$, and both refinements we tried --- an upper bound on the propensity to
consume from concavity and the sandwich, and a self-improving bootstrap on the size of the
violation --- converge to a positive bound rather than to zero.

So $\mu=3$ is not delivered by this route, and $\mu=1$ is not either, although the companion
write-up proves the latter by the rescaling argument, which is an induction on Bellman iterates
and so carries over to the life cycle unchanged. What the reduction does deliver is the end of
life, and the exact inequality that any proof must beat.

So the state of the programme is this. Theorem 1 is true in the life cycle for $\mu=3$, and the
statement to prove has the cleanest possible form, its own terminal condition propagated; the
reduction to one inequality per age is proved; and that inequality is false for the young, at
every degree of risk aversion. What is missing is an argument at zero wealth that does not go
through a pointwise comparison of the two Euler equations. The same zero-wealth state is where the
infinite-horizon slack condition was tightest, which suggests the two problems share a single
missing ingredient rather than two.

\subsection{It is not precautionary saving}

A household at zero wealth in the highest earnings state saves heavily, and the obvious diagnosis
is precautionary: it is insuring against the fall in earnings it expects. If that were the
mechanism, precautionary-saving theory would be the place to look for the missing argument. It is
not the mechanism. The same reversal happens in the \emph{certainty-equivalent} household, which
faces the conditional mean earnings path and no risk at all, and there everything is in closed
form (\texttt{numerics/precaut2.m}, \texttt{durcheck.m}).

\begingroup\small
\begin{center}
\begin{tabular}{@{}lcccc@{}}
\toprule
$\mathrm{d}(\text{saving})/\mathrm{d}r$ at zero wealth & $\mu=1$ & $\mu=3$ & $\mu=5$ & $\mu=8$\\
\midrule
Certainty equivalent (mean path, no risk) & $+21.3$ & $+2.81$ & $-0.85$ & $-2.91$\\
With risk & $+17.1$ & $+0.92$ & $-2.05$ & $-3.43$\\
\bottomrule
\end{tabular}
\end{center}
\endgroup

The sign flips between $\mu=3$ and $\mu=5$ in both columns. Bisecting, the certainty-equivalent
threshold is $\mu=4.33$ against $\mu=3.39$ with risk: risk moves the threshold by about one unit
of risk aversion and does not create the effect.

Without risk the mechanism is exact. A household with $J$ periods, no borrowing constraint and
CRRA utility has $c_{t+1}=(\beta R)^{1/\mu}c_t$ and a lifetime budget, so
\[
  c_0=\frac{W(R)}{D(R)},\qquad
  W(R)=\sum_t y_tR^{-t},\qquad
  D(R)=\sum_t\beta^{t/\mu}R^{t/\mu-t}
\]
(\lean{consumption\_mul\_pathPrice}), and saving at zero wealth rises with the rate exactly when
$W(R_2)D(R_1)\le W(R_1)D(R_2)$ (\lean{consumption\_antitone\_iff}). Both terms fall with the rate,
$W$ because future earnings are discounted harder and $D$ because the optimal path gets cheaper,
and the comparison is between the two elasticities. Differentiating,
\[
  \frac{\mathrm{d}\log c_0}{\mathrm{d}\log R}=-T_y+\Bigl(1-\frac1\mu\Bigr)T_c ,
\]
with $T_y$ and $T_c$ the Macaulay durations of the two streams, so $c_0$ falls, and saving rises,
exactly when $T_y>(1-1/\mu)T_c$. The predicted elasticity matches the exact one to three decimals
at every $\mu$ we checked. At Aiyagari's numbers $T_y=14.31$ and $T_c\approx18.6$, nearly
independent of $\mu$, so the condition is $1-1/\mu<0.769$, that is $\mu<4.33$.

The economics is then plain, and has nothing to do with risk. A higher interest rate discounts
future earnings more heavily, which makes the household poorer and pushes consumption down; it
also makes a rising consumption path cheaper, which pushes consumption today down only to the
extent that the household is willing to tilt its path. A more risk averse household will not
tilt, so the second force weakens as $\mu$ grows while the first does not, and past $\mu\approx4$
the household responds to a higher rate by consuming more today. The income stream simply has to
be longer than the consumption weights, discounted by the strength of the tilt.

\begin{theorem}[No risk, and risk aversion at most one]\label{thm:ce}
If $\mu\le1$ the comparison holds term by term, whatever the earnings profile: the
certainty-equivalent household always saves more at a higher rate.
\formal{CertaintyEquivalent.humanWealth\_mul\_pathPrice\_le\_of\_le\_one,
CertaintyEquivalent.consumption\_antitone\_of\_le\_one}
\end{theorem}

For $\mu>1$ it is the duration comparison, which is a finite computation in the primitives. So a
precautionary argument near zero wealth would be aimed at the wrong target: what has to be
controlled is the elasticity of human wealth against the elasticity of the price of the
consumption path, and risk enters only as a correction to it.

\subsection{The criterion in elasticity units, and the size of the risk correction}

The certainty-equivalent criterion is cleanest when written in the intertemporal elasticity of
substitution rather than in risk aversion. Since it is linear in $1/\mu$,
\[
  T_y>\Bigl(1-\frac1\mu\Bigr)T_c
  \qquad\Longleftrightarrow\qquad
  \frac1\mu\ >\ 1-\frac{T_y}{T_c} ,
\]
so the elasticity of substitution must exceed one minus the ratio of the two durations. At
Aiyagari's numbers that right-hand side is $0.2306$, giving $\mu<4.33$ as before, and the identity
holds to four decimals against the bisected threshold.

With risk the threshold moves down, and the move is what has to be bounded. Write
\[
  \Delta\ :=\ \frac1{\mu^{\text{risky}}}-\frac1{\mu^{\text{CE}}}
\]
for the gap in elasticity units and $s=\sigma\sqrt{1-\rho^2}$ for the conditional standard
deviation of log earnings, the dispersion the household actually faces one period ahead. Measuring
both thresholds at the same state over fifteen calibrations, $\rho\in\{0,0.6,0.9\}$ and
$\sigma$ from $0.025$ to $0.4$ (\texttt{numerics/gapscale.m}, \texttt{gapfit.m}):

\begingroup\small
\begin{center}
\begin{tabular}{@{}lcccccc@{}}
\toprule
$\rho=0.9$, $\sigma$ & $0.025$ & $0.05$ & $0.1$ & $0.2$ & $0.4$ & fitted exponent\\
\midrule
$\Delta$ & $0.0034$ & $0.0070$ & $0.0144$ & $0.0300$ & $0.0612$ & \\
$\Delta/s$ & $0.315$ & $0.320$ & $0.330$ & $0.344$ & $0.351$ & $1.04$\\
\bottomrule
\end{tabular}
\end{center}
\endgroup

The gap is first order in the conditional standard deviation, not second order: fitting
$\Delta\sim s^{b}$ gives $b=1.04$ at $\rho=0.9$, $0.89$ at $\rho=0.6$ and $0.93$ at $\rho=0$, and
the ratio $\Delta/s$ stays between $0.196$ and $0.351$ across every calibration we ran. That
first-order behaviour is itself informative: a precautionary effect on saving would be second
order in the dispersion, whereas the probability that the borrowing constraint binds in a
low-earnings state next period moves at first order, and that is our reading of where the
correction comes from.

Taking the largest ratio observed as the bound, $\Delta\le0.36\,s$, gives a criterion entirely in
the primitives:
\begin{equation}\label{eq:practical}
  \frac1\mu\ \ge\ 1-\frac{T_y}{T_c}\ +\ 0.36\,\sigma\sqrt{1-\rho^2}.
\end{equation}
At Aiyagari's calibration the right-hand side is $0.2306+0.0628=0.2934$, that is $\mu\le3.41$,
against the $3.39$ that bisection on the full life-cycle model gives. The bound therefore
reproduces the measured threshold to within two per cent, and it says the same thing the earlier
sections said: the calibration he used for $\mu=3$ is inside, and $\mu=5$ is not.

\subsubsection*{The certainty-equivalent half, for every risk aversion}

The first inequality of \eqref{eq:practical} was a theorem only for $\mu\le1$, where the exponent
$\theta=1-1/\mu$ is nonpositive and every factor moves the same way. It is now a theorem for every
$\mu$. Writing $\lambda=R_1/R_2<1$, the criterion $W(R_2)D(R_1)\le W(R_1)D(R_2)$ compares how the
same weights respond to $\lambda^t$ against $(\lambda^t)^\theta$. Each $\lambda^t$ lies in
$[\lambda^{J-1},1]$, and on that interval the concave $x\mapsto x^\theta$ lies above its chord,
which is affine; replacing the power by its chord leaves a comparison of two sums linear in
$\lambda^t$.

\begin{theorem}[The chord criterion]\label{thm:chord}
For $\mu>1$, if
$W(R_2)D(R_1)\le W(R_1)\bigl(\iota\,D(R_1)+\varsigma\,D_\lambda(R_1)\bigr)$, where $\iota$ and
$\varsigma$ are the intercept and slope of that chord and $D_\lambda$ is the path price with each
term discounted once more by $\lambda^t$, then initial consumption falls with the rate.
\formal{CertaintyEquivalent.consumption\_antitone\_of\_chord},
\lean{CertaintyEquivalent.chord\_le\_pathPrice}
\end{theorem}

The replacement costs almost nothing, because it is first-order exact: as $R_2\to R_1$ it
reproduces $T_y\ge(1-1/\mu)T_c$ exactly. At Aiyagari's numbers the chord threshold is $4.3263$
against the exact $4.3267$, and the cost grows only with the size of the rate step, to $0.04$ at
$\Delta r=10^{-3}$ and $0.40$ at $10^{-2}$ (\texttt{numerics/chord.m}). Together with the $\mu\le1$
theorem this settles the certainty-equivalent household at every risk aversion.

\subsubsection*{The risk half, and why it resists}

Two caveats remain. The constant $0.36$ is measured over fifteen calibrations and
is not proved. And both are statements at zero wealth in the highest earnings state, which is
where the life-cycle model binds; the reduction of Theorem~\ref{thm:reduction} is what would carry
them to the rest of the state space, and its own condition fails, so this is a sharp description
of where the truth lies rather than a proof of it.

It is worth being precise about why the constant resists. What has to be bounded is a difference of
\emph{rate responses} between two problems, the risky one and its certainty equivalent, and the
threshold is the rate at which one of those responses changes sign. No bound on the \emph{level} of
consumption helps with that, however sharp: at Aiyagari's calibration the sharpest sandwich
available has width $0.30$ to $0.48$ at the state in question, while the response itself runs from
$-2.48$ at $r=2\%$ through zero to $+0.29$ at $r=6\%$ (\texttt{numerics/scale.m}). A bound on where
consumption lies says nothing about how it moves. What is needed is a bound on the rate response
itself, and that is exactly the object the companion write-up proved cannot be had from the affine
induction, whose per-step multiplier is at least $1/\Thorn>1$. The two halves of
\eqref{eq:practical} therefore stand on different ground: the first is now a theorem for every
risk aversion, and the second is blocked by the same obstruction that closed the rate-response
bound.

\subsubsection*{How much slack a bound on the correction has}

That said, ``blocked'' understates how good the position is, and the quantity worth measuring is
not the correction but the room around it. Write the risky household's saving response at the state
as the certainty-equivalent response less a correction,
$\partial g/\partial r=\partial g_{\mathrm{CE}}/\partial r-\mathrm{corr}$. Theorem 1 at the state
follows from any bound $B\ge\mathrm{corr}$ with $B\le\partial g_{\mathrm{CE}}/\partial r$
(\lean{CertaintyEquivalent.saving\_nonneg\_of\_correction\_le}), so what matters is the ratio
$(\partial g_{\mathrm{CE}}/\partial r)/\mathrm{corr}$: a bound loose by less than that factor still
delivers the theorem. Measured at $r=4\%$, zero wealth, top earnings state:

\begin{center}
\begin{tabular}{lrrrrrrr}
\toprule
$\gamma$ & $1$ & $2$ & $2.5$ & $3$ & $3.2$ & $3.39$ & $5$\\
\midrule
$\partial g_{\mathrm{CE}}/\partial r$ & $20.98$ & $7.28$ & $4.55$ & $2.73$ & $2.17$ & $1.69$ & $-0.89$\\
$\partial g/\partial r$ (risky) & $17.08$ & $4.88$ & $2.49$ & $0.92$ & $0.44$ & $0.03$ & $-2.05$\\
correction & $3.89$ & $2.40$ & $2.06$ & $1.82$ & $1.73$ & $1.66$ & $1.15$\\
\midrule
slack factor & $\mathbf{5.4}$ & $\mathbf{3.0}$ & $\mathbf{2.2}$ & $\mathbf{1.5}$ & $1.25$ & $1.02$ & ---\\
\bottomrule
\end{tabular}
\end{center}

Two things stand out. The correction is \emph{stable}: it stays between $1.15$ and $3.89$ across
$\gamma\in[1,5]$, and it \emph{falls} as risk aversion rises, while the certainty-equivalent response
collapses from $20.98$ to below zero. The threshold is therefore set by the certainty-equivalent side
and not by risk at all --- which is the same conclusion the closed form gave, now with the correction
measured rather than inferred. And the slack is generous where it matters: a bound on the correction
may be wrong by a factor of $5.4$ at $\gamma=1$, $3.0$ at $\gamma=2$ and $1.5$ at $\gamma=3$ and
still certify Theorem 1 there.

That is a far weaker demand than the one left open in Section~\ref{sec:inc}, where the remaining
estimate needs an absolute error below $0.03$ on a quantity of size about one. Here a crude
precautionary bound would do. Note only that the slack is \emph{not} uniform in $\gamma$: the
correction peaks at $3.89$ at $\gamma=1$ while $\partial g_{\mathrm{CE}}/\partial r$ is only $2.73$
at $\gamma=3$, so a single constant cannot cover the range --- the bound must be allowed to depend on
$\gamma$, which costs nothing, since Theorem 1 is wanted at a calibration and not as a family.
Numerics: \texttt{riskgap.m}.

\subsubsection*{The precautionary split, and what a bound on it must achieve}

The slack above is measured on derivatives at $r=4\%$. Theorem 1 is a statement about an interval, so
it is worth redoing in secant form, where the object to be bounded becomes concrete. Write the risky
household's saving at the state as the certainty-equivalent saving plus precautionary saving,
$g=g_{\mathrm{CE}}+P$ with $P=c_0^{\mathrm{CE}}-c_0\ge0$. Then Theorem 1 across $[r_1,r_2]$ is
\emph{exactly}
\begin{equation}\label{eq:precfall}
  P(R_1)-P(R_2)\;\le\;g_{\mathrm{CE}}(R_2)-g_{\mathrm{CE}}(R_1)\;=:\;\text{margin},
  \tag{P}
\end{equation}
a bound on how far precautionary saving falls across the interval
(\lean{CertaintyEquivalent.saving\_le\_of\_precautionary\_fall}). No derivative appears. Two weakenings
of \eqref{eq:precfall} are worth separating, and they close at different risk aversions. Measured over
$r\in[4\%,6\%]$ at zero wealth in the top earnings state:

\begin{center}
\begin{tabular}{lrrrrrr}
\toprule
$\gamma$ & $1$ & $2$ & $2.5$ & $3$ & $3.2$ & $3.39$\\
\midrule
$P(R_1)$ & $0.109$ & $0.154$ & $0.176$ & $0.198$ & $0.206$ & $0.214$\\
$P(R_1)-P(R_2)$ & $0.044$ & $0.032$ & $0.028$ & $0.025$ & $0.024$ & $0.023$\\
margin & $0.329$ & $0.107$ & $0.062$ & $0.031$ & $0.022$ & $0.014$\\
\midrule
level route: $P(R_1)\le$ margin & $\checkmark$ & $\times$ & $\times$ & $\times$ & $\times$ & $\times$\\
$\theta$ admissible & --- & $0.305$ & $0.650$ & $0.842$ & $0.895$ & $0.936$\\
$\theta$ true & $0.598$ & $0.793$ & $0.839$ & $0.872$ & $0.883$ & $0.892$\\
\bottomrule
\end{tabular}
\end{center}

At $\gamma=1$ the crudest weakening works: discard $P(R_2)\ge0$ and a bound on the \emph{level} of
precautionary saving suffices, with the margin $0.329$ exceeding $P(R_1)=0.109$ threefold
(\lean{saving\_le\_of\_precautionary\_level}). Above that the level route dies and one needs the
multiplicative form: keeping a fraction $\theta$ of precautionary saving at the higher rate,
\eqref{eq:precfall} becomes $(1-\theta)P(R_1)\le\text{margin}$
(\lean{saving\_le\_of\_precautionary\_ratio}). At $\gamma=2$ the admissible $\theta$ is $0.305$ against
a true ratio of $0.793$: precautionary saving need only be shown not to lose seventy per cent of
itself when the rate rises by two points, which is a crude statement. At $\gamma=3$ the admissible
$\theta$ is $0.842$ against $0.872$, a margin of three per cent. Above $\gamma\approx3.1$ there is no
margin at all, because Theorem 1 itself fails on that interval --- the $3.39$ of the derivative
comparison is the threshold at $r=4\%$, not across $[4\%,6\%]$.

So the honest target is $\gamma\le2$, and what it needs is two-sided bounds on precautionary saving at
one state accurate to about a quarter: measured, $P(R_1)=0.154$ and $P(R_2)=0.122$ against a required
gap of $0.107$. The sandwich available gives $\kappa_kH_k=1.40$ at that state, nine times too loose,
so this is a new estimate rather than an assembly of existing ones --- but it is a bound on a level,
at one state, with a factor of four to give away, which is a different kind of demand from anything
the increment route leaves open. Numerics: \texttt{precbound.m}.

\subsection{The channel is the borrowing constraint}

The first-order scaling points at the constraint rather than at precaution, and the point can be
tested directly: loosen the borrowing limit and watch the threshold. If the constraint is the
channel, the risky threshold should walk back to the certainty-equivalent one as the limit
approaches the natural one, $y_{\min}/r$. At Aiyagari's calibration it does
(\texttt{numerics/borrow.m}), and it does so almost linearly:

\begingroup\small
\begin{center}
\begin{tabular}{@{}lcccccccc@{}}
\toprule
Borrowing limit & $0$ & $0.5$ & $1$ & $2$ & $3$ & $4$ & $5$ & $6$\\
\midrule
Threshold $\mu$ & $3.42$ & $3.50$ & $3.58$ & $3.74$ & $3.89$ & $4.05$ & $4.21$ & $4.38$\\
Share of the gap closed & $0\%$ & $11\%$ & $21\%$ & $40\%$ & $57\%$ & $74\%$ & $90\%$ & $104\%$\\
\bottomrule
\end{tabular}
\end{center}
\endgroup

The natural limit here is $y_{\min}/r=6.75$, and by a limit of $6$ the gap is gone: the risky
household's threshold, $4.38$, is the certainty-equivalent household's, $4.33$. So at this
calibration the borrowing constraint accounts for the whole of the correction, and precaution
accounts for none of it. Precaution is certainly present --- at the loose limit the risky
household still saves $1.62$ against the certainty-equivalent $1.46$, a precautionary excess of
eleven per cent --- but it moves the level of saving without moving the sign of its response to
the interest rate. Across the other calibrations we ran (\texttt{numerics/borrow3.m}) loosening
the limit always closes the gap and sometimes overshoots it, so the constraint is the dominant
channel rather than an exact decomposition: with ample borrowing, risk can make Theorem 1 easier
than certainty equivalence rather than harder.

For the programme this is the useful part. The missing ingredient for $\mu=3$ is not a
precautionary bound but a bound on what the borrowing constraint does to the comparative static,
and that is a more tractable object: without a binding constraint the household's Euler equation
holds with equality at every state and Theorem 1 at zero wealth reduces to the duration criterion,
which is proved outright for $\mu\le1$ and is an explicit finite computation in the primitives
otherwise. What remains is the constrained states, which at this calibration are about eight per
cent of the state space at every age.

\subsection{The constraint channel is monotone, and needs no condition}

The comparative static in the borrowing limit itself is a theorem, and unlike the one in the
interest rate it needs nothing beyond the standing assumptions.

\begin{theorem}[A tighter limit means less consumption and more saving]\label{thm:limit}
Two households share every primitive and differ only in the borrowing limit, $f_2\le f_1$. Then at
every age and every state in the tighter household's region it consumes less and saves more.
\formal{IncomeFluctuation.finiteConsumption\_le\_of\_floor\_le,
IncomeFluctuation.le\_finitePolicy\_of\_floor\_le}
\end{theorem}

The induction is the one of Theorem~\ref{thm:reduction} with the interest rate held fixed, and
that is exactly what removes the difficulty. Suppose the claim holds with $k$ periods to come and
fails at $k+1$, so the tighter household consumes strictly more and therefore saves strictly less.
Its Euler inequality under the cap, the looser household's Euler inequality above its floor ---
available because the looser household saves strictly more than the tighter one and hence strictly
more than either floor --- the induction hypothesis, and monotonicity of consumption in wealth
chain into $u'(c_1)\ge u'(c_2)$, contradicting $c_1>c_2$. The gross return appears on both sides
and cancels, so there is no substitution term to control and no analogue of
\eqref{eq:stageslack} is needed. Tightening the limit raises the marginal value of saving
unambiguously, and the terminal stage is again free: a household with a lower floor runs down
further in its last period.

That is the one-sided half of what the numbers showed, and it says the constraint channel is
better behaved than the rate channel rather than merely different. What it does not yet give is a
\emph{size}: Theorem~\ref{thm:limit} orders the two consumption functions without bounding the
distance between them, and it is the distance that turns into the $0.9$ units of risk aversion
between the constrained and certainty-equivalent thresholds. The natural next object is therefore
a modulus for Theorem~\ref{thm:limit}, and it has a clean form.

\begin{theorem}[The modulus: a relaxation of $\Delta$ is worth at most $\Delta$ of wealth]
\label{thm:modulus}
Let $\Delta=f_1-f_2\ge0$ and $r\ge0$. Then at every age and every state,
\[
  c^{f_1}_k(x)\ \le\ c^{f_2}_k(x)\ \le\ c^{f_1}_k(x+\Delta) .
\]
\formal{IncomeFluctuation.finiteConsumption\_le\_of\_floor\_le,
IncomeFluctuation.finiteConsumption\_le\_shift}
\end{theorem}

The household that may borrow down to $f_2$ therefore behaves, at wealth $x$, no more freely than
the household confined to $f_1$ would at wealth $x+\Delta$: the entire value of the extra
borrowing capacity is bounded by the capacity itself. The proof is the same backward induction,
carried out at shifted arguments, and the shift pays for itself twice. The terminal stage holds
with slack, since the looser household gains $\Delta$ by running down further while the shifted
comparison gives it $(1+r)\Delta$; and the interiority the step needs comes free, because the
failure hypothesis forces the tighter household's saving to exceed the looser one's by more than
$\Delta$, hence to exceed $f_2+\Delta=f_1$.

Two consequences. Since saving never falls with wealth, the sandwich gives the plain bound
$0\le c^{f_2}_k(x)-c^{f_1}_k(x)\le(1+r)\Delta$, so the effect of the limit on consumption is at
most the relaxation, up to one period's interest. And because the bound is stated as a shift in
wealth rather than as an additive constant, it composes with the minimal propensity to consume:
whatever upper bound on the propensity to consume out of wealth one has at age $k$, the modulus
turns it directly into a bound on the consumption gap, and it is that product which becomes the
gap between the constrained and certainty-equivalent thresholds. \subsection{The upper bound on the propensity to consume, and what it is worth}

The missing piece is an upper bound on the propensity to consume out of wealth. The repository has
the minimal propensity, which is a lower bound; the upper one follows by the same two-point Euler
comparison, with no derivative taken. Write the bound as a secant slope in assets, so that $s_k$
bounds $c_k(a_2)-c_k(a_1)$ by $s_k(a_2-a_1)$.

\begin{theorem}[The propensity to consume falls with the horizon]\label{thm:mpc}
In the last period the slope is exactly $R$ (\lean{stageConsumption\_sub\_zero}). If at stage $k$
the slope is at most $s$ and consumption is at least $c_{\min}$, then at a state where the poorer
household saves, the slope at stage $k+1$ is at most
\[
  R\cdot\frac{\nu}{1+\nu},\qquad \nu=\frac{c\,s}{c_{\min}},
\]
with $c$ the poorer household's consumption.
\formal{IncomeFluctuation.stageConsumption\_sub\_le\_of\_slope}
\end{theorem}

The step is the two Euler inequalities at the two wealth levels: the induction hypothesis bounds
the rise in tomorrow's consumption by $s\delta$ with $\delta$ the difference in saving, which
bounds the fall in marginal utility by the factor $(1+s\delta/c_{\min})^{-\mu}$, and the budget
identity $c_2-c_1=R(a_2-a_1)-\delta$ turns that into the stated slope. Since $\nu\mapsto\nu/(1+\nu)$
is a contraction towards $1-1/\nu$, the slope falls from $R$ towards $R(1-c_{\min}/(c\,s))$ as the
horizon lengthens: the longer the life, the less of an extra unit of wealth is spent at once,
which is the qualitative fact one wants.

It is not sharp enough. At Aiyagari's calibration, read at zero wealth in the highest earnings
state, consumption is about $1.32$ and the worst consumption next period about $0.88$, so
$\nu\approx1.5\,s$ and the recursion settles at a propensity of about $0.36$, against a true
propensity there of about $0.04$. Fed through Theorem~\ref{thm:modulus} it bounds the effect of
the borrowing limit on consumption by about $2.5$ where the truth is about $0.04$, a factor of
sixty. The loss is the ratio of consumption today to the worst consumption tomorrow, and that
ratio is where the sandwich's slack of order human wealth enters --- the same obstruction that
closed the infinite-horizon routes in the companion write-up, met here for the fourth time and in
the same form.

So the chain now has every link proved except its sharpness. Theorem~\ref{thm:ce} gives the
criterion without risk, Theorem~\ref{thm:limit} signs the constraint's effect,
Theorem~\ref{thm:modulus} bounds it by a shift in wealth, and Theorem~\ref{thm:mpc} converts a
shift in wealth into consumption. What none of them supplies is a bound on consumption that is
tight to a few per cent rather than to a factor of ten, and every route we have tried reduces to
exactly that.

\subsection{Carrying earnings across all remaining ages}

The two halves of the sandwich are not equally to blame. The upper half,
$c_k\le\kappa_k(m+H_k)$, carries the expected present value of future earnings and is about twenty
per cent above the truth at Aiyagari's calibration. The lower half, $\kappa_k m\le c_k$, ignores
future earnings entirely and is a factor of ten below it. The repair is the same device applied on
the other side: carry earnings forward across all the remaining ages instead of one step at a
time.

\begin{theorem}[The lower half, with human wealth]\label{thm:lower}
Let $H_k$ be the present value of $k$ periods of the minimum earnings, $H_0=0$ and
$RH_{k+1}=y_{\min}+H_k$. Then $\lambda_k(m+H_k)\le c_k(m)$ at every age and state, where
$\lambda_0=1$ and $\lambda_{k+1}=\Thorn\lambda_kR/(\Thorn+\lambda_kR)$.
\formal{IncomeFluctuation.minHumanWealth\_le\_stageConsumption}
\end{theorem}

The old bound is the case $H=0$. The propensity $\lambda_k$ is the finite-horizon propensity of
Theorem~\ref{thm:sandwich} discounted once more by patience at each step, and the extra discount is
not decoration: without it the bound is false at the borrowing constraint, where the household
consumes exactly its cash on hand while $\kappa_k(m+H_k)$ exceeds it, because $H_k$ discounts
earnings at $R$ and $1/\kappa_k$ discounts at $R/\Thorn$. Discounting $H_k$ at $R/\Thorn$ instead
repairs the corner and breaks the induction; the extra patience factor repairs both at once
(\lean{lowMPC\_mul\_minHumanWealth\_le}), and at $\beta=0.96$, $r=4\%$ and $\mu=3$ it costs about
three per cent over sixty periods.

At the state where the life cycle binds the household enters the next period with cash on hand
about $2.0$ and consumes about $0.88$. The old bound gives $0.086$; this one gives $0.33$. Through
Theorem~\ref{thm:mpc} that improves the bound on the propensity to consume from about $0.94$ to
about $0.75$, against a true propensity of $0.04$.

So the multi-age argument buys a factor of four and not a factor of ten, and the reason is visible
in the statement: it carries the \emph{minimum} earnings, and a household that drew the minimum
forever really would consume this little. What the programme needs is a lower bound carrying the
\emph{expected} human wealth, and that is not available from any chain of Euler inequalities,
because the Euler inequality at a constrained state is exactly the statement that expected future
earnings cannot be borrowed against. That is the one quantitative fact about a single household
on which the whole programme turns, and identifying it precisely is what makes it attackable.

\subsection{The constraint-incidence bound}

Both halves can be carried at once. Accumulate human wealth by the \emph{risk-adjusted}
expectation --- the power mean with exponent $-\mu$, which is what the Euler equation itself
produces --- and then \emph{cap} it state by state at the level the borrowing constraint permits:
\[
  H_0=0,\qquad
  H_{k+1}(z)=\min\Bigl\{\tfrac1R\,M^z_{-\mu}\bigl(y_w+H_k(w)\bigr),\
  \Bigl(\tfrac1{\kappa_{k+1}}-1\Bigr)y_z\Bigr\}
\]
(\lean{riskHumanWealth}), with $\kappa$ the propensity of Theorem~\ref{thm:sandwich}.

\begin{theorem}[The constraint-incidence bound]\label{thm:incidence}
$\kappa_k\bigl(m+H_k(z)\bigr)\le c_k(a,z)$ at every age and state.
\formal{IncomeFluctuation.riskHumanWealth\_le\_stageConsumption}
\end{theorem}

The two ingredients meet exactly where they should. The power mean is what the Euler inequality
delivers, and the step needs only that it is superadditive under a shift, which follows from its
concavity and homogeneity (\lean{powerMean\_add\_const\_le}). The cap is the constraint's
incidence: at a corner the household consumes its cash on hand, which is at least $y_z$, and the
cap is precisely the largest human wealth the bound may claim without exceeding that. No extra
patience factor is needed, because the cap does the work the patience factor did in
Theorem~\ref{thm:lower}.

At Aiyagari's calibration, in the first period of a sixty-period life
(\texttt{numerics/incidence.m}), the cap binds at the two lowest earnings states and nowhere else,
which is where a household at zero assets is against the limit. The bound reaches

\begingroup\small
\begin{center}
\begin{tabular}{@{}lccc@{}}
\toprule
At zero wealth & lowest earnings & middle & highest\\
\midrule
Bound as a fraction of the truth & $1.000$ & $0.890$ & $0.877$\\
Minimum-earnings bound, for comparison & $1.010$ & $0.409$ & $0.295$\\
\bottomrule
\end{tabular}
\end{center}
\endgroup

and over the whole grid it never exceeds the truth, touching it exactly at the constrained corner.
So the consumption bound is now within about twelve per cent where the life cycle binds, against a
factor of three before. That is the sharpest bound in the development, and it was reached by
bounding where the constraint binds rather than by another Euler chain.

It does not improve the propensity bound of Theorem~\ref{thm:mpc}, because that bound is driven by
the worst earnings state next period, which is exactly where the cap binds. Closing the threshold
argument would need the same sharpening on the worst state, and that turns out to be impossible
within this class.

\subsection{The worst state is pinned}

Two facts settle it. The first is that the cap is not conservatism but necessity.

\begin{theorem}[The cap is necessary]\label{thm:cap}
Wherever a household at zero assets is against the borrowing limit, every affine lower bound
$\kappa_k(m+H)$ valid at that state obeys $H\le(1/\kappa_k-1)y_z$.
\formal{IncomeFluctuation.cap\_necessary\_of\_corner}
\end{theorem}

The proof is one line --- at the corner consumption is exactly $y_z$ --- and the content is that
Theorem~\ref{thm:incidence} claims the largest human wealth any bound of its shape may claim. It
cannot be sharpened within the affine class at all.

The second is that dropping the cap does not merely weaken the proof, it makes the bound false.
Measuring effective human wealth as $c/\kappa-m$ at next-period assets of $1.67$
(\texttt{numerics/worst.m}):

\begingroup\small
\begin{center}
\begin{tabular}{@{}lcccc@{}}
\toprule
Effective human wealth at & truth & capped mean & uncapped mean & expectation\\
\midrule
lowest earnings & $10.96$ & $5.99$ & $13.11$ & $16.57$\\
highest earnings & $27.92$ & $23.90$ & $26.38$ & $32.24$\\
\bottomrule
\end{tabular}
\end{center}
\endgroup

At the lowest earnings state the uncapped power mean, $13.11$, \emph{exceeds} the truth, $10.96$,
so a bound built on it is invalid: the household will be against the limit later, and the uncapped
recursion cannot see that. The truth sits between the two, and the distance is the price of
capping at zero assets when the household in fact holds some. A bound that closed it would have to
let human wealth depend on the assets held, which requires a lower bound on the household's own
asset path; the one available, from the upper sandwich, is negative at the state that matters.

So the worst state is pinned from both sides: the cap cannot be raised, because the corner forbids
it, and it cannot be dropped, because the truth is below what dropping it would claim. That closes
this family of arguments, and with it, for now, the route from a consumption bound to the
threshold. What the programme has instead is an exact statement of what a sharper argument must
supply: a lower bound on consumption that is asset-dependent rather than affine, and the asset
path to support it.

\subsection{The bound with assets in it}\label{sec:assetdep}

The obstruction just described is no longer binding, because the path bound it asked for is
available. Proposition~\ref{prop:floor} read at the household's own earnings state rather than at
the worst one is positive exactly where the cap bites hardest: a household in a high earnings
state saves, so a period later it is not at the corner, and the cap that then applies to it is the
one for a household with assets.

So assets go through the recursion twice over. Writing $\Phi_k(z,a)$ for that saving floor,
clamped at zero,
\[
  G_0(z,a)=0,\qquad
  G_{k+1}(z,a)=\min\Bigl(\tfrac1R M^z_{-\mu}\bigl(y_w+G_k(w,\Phi_{k+1}(z,a))\bigr),\;
    \bigl(\tfrac1{\kappa_{k+1}}-1\bigr)(y_z+Ra)\Bigr),
\]
and the bound is $\kappa_k(m+G_k(z,a))\le c_k(a,z)$. The cap now grows with assets, and tomorrow's
human wealth is read at the assets the household will hold rather than at zero.

\begin{theorem}[The asset-dependent incidence bound]\label{thm:assetdep}
$\kappa_k\bigl(m+G_k(z,a)\bigr)\le c_k(a,z)$, and $G_k(z,a)\ge H_k(z)$ for $a\ge0$, so the bound is
never worse than Theorem~\ref{thm:incidence}.
\formal{IncomeFluctuation.assetHumanWealth\_le\_stageConsumption},
\lean{IncomeFluctuation.riskHumanWealth\_le\_assetHumanWealth}
\end{theorem}

Measured against the truth over the whole grid at three ages
(\texttt{numerics/assetdep.m}, \texttt{numerics/worstgrid.m}):

\begingroup\small
\begin{center}
\begin{tabular}{@{}lccc@{}}
\toprule
Worst ratio of bound to truth & $r=-7\%$ & $r=4\%$ & $r=10\%$\\
\midrule
Theorem~\ref{thm:incidence}, cap at zero assets & $0.237$ & $0.606$ & $0.760$\\
Theorem~\ref{thm:assetdep}, cap at the assets held & $0.350$ & $0.746$ & $0.857$\\
\bottomrule
\end{tabular}
\end{center}
\endgroup

The new bound never exceeds the truth anywhere on the grid. The gain is largest exactly where the
old one was weakest, at positive assets in low earnings states: at $r=4\%$ with five units of
assets and the lowest earnings, the ratio rises from $0.621$ to $0.835$. Where the old bound was
already good, at zero assets, little changes, which is as it should be, since that is the state
the cap was computed for.

What the refinement does not improve is the ceiling on aggregate capital, because the state that
binds there is the highest earnings one, where the cap was never active: the implied ceiling at
$r=-7\%$ is $22.15$ against $22.17$. That is still far below the $39.66$ that feasibility alone
gives, which is worth carrying into the existence argument of Section~\ref{sec:eqm}.

\section{Theorem 1 at logarithmic utility}\label{sec:log}

The slack condition of Section~\ref{sec:eqm} is an Euler condition, and it is worth asking whether
it settles the one case the stationary model does settle, $\mu=1$. It does not, and the reason is
instructive. Checking it on the whole grid at $\gamma=1$ (\texttt{numerics/slacklog.m}), it holds
at $r=0$ and $r=2\%$ and fails at $r=4\%$ and above, by about twelve per cent in elasticity terms,
at the youngest age with zero assets in the highest earnings state. Since a life-cycle equilibrium
sits above the rate of time preference, $4.17\%$ here, the Euler reduction says nothing where it
is wanted, even at logarithmic utility.

\subsection{A condition on the value gap}

A different sufficient condition does hold. Write $V_k(\cdot;R)$ for the stage-$k$ value.

\begin{definition}[Value gap]\label{def:gap}
The value gap has slope at least $\theta$ at stage $k$ if, for $x\le y$ in $[0,\bar a]$,
\[
  \theta\,(y-x)\ \le\ \bigl[V_k(y;R_2)-V_k(y;R_1)\bigr]-\bigl[V_k(x;R_2)-V_k(x;R_1)\bigr].
\]
\formal{IncomeFluctuation.ValueGap}
\end{definition}

By the envelope theorem the slope of the gap is $R_2u'(c_2)-R_1u'(c_1)$, so the condition says the
marginal value of assets rises with the rate. At $\gamma=1$ that reads $c(R_2)/R_2\le c(R_1)/R_1$:
consumption rises by proportionally less than the rate does.

\begin{theorem}[Theorem 1 from the value gap]\label{thm:gapmono}
If the value gap has a positive slope at stage $k$, then saving at stage $k$ rises with the rate,
at every state. Aggregate capital rises with it.
\formal{IncomeFluctuation.stagePolicy\_mono\_of\_valueGap},
\lean{IncomeFluctuation.olgCapital\_le\_of\_valueGap}
\end{theorem}

Unlike the Euler route this needs no induction over ages. The condition on the continuation gives
the conclusion for the household that faces it, directly, and the argument is Topkis with two
continuations: the two households differ only in their cash on hand, which is ordered by the rate,
and in their continuations, whose difference is increasing, so the higher-rate household cannot
save less. Only concavity of period utility is used, as in Section~\ref{sec:ret}.

\subsection{Why it holds at $\mu=1$, and only near there}

At $\gamma=1$ the patience factor is $\text{\th}=\beta R$, so $q=\text{\th}/R=\beta$ and the
propensity $\kappa_k=1/\sum_{i\le k}\beta^i$ of Proposition~\ref{prop:kappa} does not depend on
the interest rate at all. Human wealth does, and falls as the rate rises. Where the borrowing
constraint is slack, $c=\kappa_k(\bar y+Ra+H)$ and
\[
  R_2c(R_1)-R_1c(R_2)\ =\ \kappa_k\bigl[(R_2-R_1)\bar y+R_2H_1-R_1H_2\bigr]\ >\ 0 ,
\]
which is the condition, with room. At a corner, $c=\bar y+Ra$ and the same inequality reduces to
$R_1\bar y\le R_2\bar y$.

Numerically (\texttt{numerics/condB.m}) the condition holds at $\gamma=1$ at every age, state and
rate tested, and fails from $\gamma=2$, where between $19\%$ and $51\%$ of states violate it. It is
therefore sufficient and not necessary: Theorem 1 itself still holds at $\gamma=2$, and up to about
$3.4$. So this route covers $\mu=1$ and closes nothing above it.

What is proved, then, is the reduction, and it is a reduction to a condition that is true at
logarithmic utility rather than one that is false there. Discharging the condition itself from
primitives is not available from the sandwich of Theorem~\ref{thm:sandwich}: the sandwich's width
is $\kappa_kH_k$, and the rate response of exactly that quantity is what drives the gap, so the
bounds cannot see it. That is the same obstruction the consumption-bound family ran into above,
and it is the piece a sharper argument would have to supply.
